\documentclass[11pt]{article}
\usepackage[margin=1in]{geometry}
\usepackage{amsmath,amssymb,amsthm,booktabs,array,graphicx,xcolor}
\usepackage[hidelinks]{hyperref}
\hypersetup{pdftitle={Compositional Adaptation with Finite-Gallery Retrieval Constraints},pdfauthor={Private working manuscript},pdfsubject={Audited V10 seed-17 development screens pass; replication and practical evaluation incomplete}}
\setlength{\emergencystretch}{2em}
\newtheorem{proposition}{Proposition}
\newcommand{\ip}[2]{\langle #1,#2\rangle}
\newcommand{\norm}[1]{\lVert #1\rVert}
\newcommand{\R}{\mathbb R}
\newcommand{\indicator}{\mathbf 1}
\graphicspath{{results/practical_v9/figures/}{../results/practical_v9/figures/}{results/practical_v10/figures/}{../results/practical_v10/figures/}}
\title{Compositional Adaptation with Finite-Gallery Retrieval Constraints}
\author{Private working manuscript}
\date{5 October 2026}
\makeatletter
\def\ps@draftstatus{%
  \def\@oddhead{{\slshape Private working draft --- V10 in progress}\hfil\thepage}%
  \let\@evenhead\@oddhead
  \def\@oddfoot{}\let\@evenfoot\@oddfoot}
\makeatother
\begin{document}
\pagestyle{draftstatus}
\maketitle
\begin{center}
\fbox{\parbox{.91\linewidth}{\textbf{Experiment status.}
This is an incomplete research manuscript, not a submission.
The V10 seed-17 development screens pass for both encoders and have been
independently checked. Replication and matched controls remain incomplete;
no V10 practical benchmark result is reported. Earlier V8 and V9 pilots
failed their development gates. The practical goal has not been established.}}
\end{center}

\begin{abstract}
Compositional adaptation of image--text models should be evaluated together
with the retrieval behavior it changes. We study this problem using a
single score on frozen embeddings. A centered bilinear residual is trained
with full-gallery source retrieval and a joint loss that ranks source and
image-supported captions above contradictions. Explicit constraints retain
a fraction of each protected frozen retrieval margin in both directions.
For fixed galleries, we prove preservation of every previously correct
top-1 query, including deterministic ties, and describe a numerical
certificate that checks the complete gallery. The objective is convex,
although our finite optimizer is not claimed to find its exact optimum.
Independently audited fits preserve all protected training queries for
ViT-B/32 LAION2B and RN50 OpenAI. Two pilots on 1,200 owners fail their
development gates despite training retention. An expanded 6,000-owner
recipe passes both seed-17 development screens: image-to-text R@1 changes
by $+1.00$ and $+1.44$ percentage points, with adjusted lower bounds of
$-.78$ and $-.89$ above the declared $-1$-point tolerance. The respective
source/support proxy gains are $+.67$ and $+1.08$ points, but both gain
intervals include zero. The expansion also increases optimizer updates,
so its effect cannot be attributed to data alone. Replication, matched
controls and practical benchmark evaluation remain incomplete. The
supported result is a finite-gallery certificate and a passed initial
development screen, not a demonstrated practical solution.
\end{abstract}

\section{Introduction}
Contrastive image--text models such as CLIP~\cite{clip} support retrieval
through a shared embedding space. Adapting their scoring geometry can
improve discrimination between closely related captions, but retrieval
places a different demand on the same score: each query must compete with
an entire gallery. A model can distinguish an annotated contradiction more
reliably while changing which unrelated image or caption ranks first.
Evaluating only a small set of candidate captions therefore does not
establish that retrieval has been retained.

The relation between positive captions also matters. A source caption
identifies a paired image in a retrieval dataset. Another description may
be supported by that image without being a paraphrase of the source. A
contradiction asserts incompatible content, whereas a neutral description
is not necessarily false. Treating these relations as interchangeable
would alter both supervision and the meaning of the evaluation. In
particular, a source/support ranking test is not the lexical-invariance
test in SugarCrepe++, whose positive descriptions are intended to be
semantically equivalent~\cite{sugarcrepepp}.

We formulate adaptation with these distinctions explicit. Source ownership
defines a full-gallery bidirectional retrieval objective. A separate
relation objective requires a source caption and an image-supported
caption to outrank a contradiction. The two objectives share one learned
score. Rather than asking a retention penalty alone to preserve existing
decisions, we constrain every previously correct training query against
every irrelevant training-gallery competitor. Incorrect frozen decisions
remain free to improve.

The score is a centered, low-dimensional bilinear residual on fixed
embeddings. Its parameters enter scores affinely, yielding a convex
objective and linear margin constraints inside a norm ball. This structure
makes a finite-gallery retention proposition possible and permits direct
numerical verification. It does not create an unseen-data guarantee.
The empirical question is whether the constraint mechanism improves the
composition--retrieval tradeoff outside the fitting gallery.

Our contribution is narrowly scoped to this constrained formulation, its
tie-aware finite-gallery analysis, and an auditable evaluation program
separating training retention from transfer. Linear alignment is already
studied by LABCLIP~\cite{labclip}; compositional improvement with multimodal
retention is already studied by CLIC and FSC-CLIP~\cite{clic,fsc}.
We make no priority claim for those broader ideas. Nor do we presently
claim an empirical advance over those methods.
Protecting earlier correct predictions is also established in model-update
research~\cite{pct,rcat}; our scope is its explicit realization for this
fixed-feature, bidirectional full-gallery problem.

The completed experiments provide a concrete warning against inferring
transfer from a training certificate. Both V8 and V9 preserve their
protected training queries. V8 fails development; V9 improves all four
development means but still fails the adjusted image-to-text retention
criterion. An expanded experiment, V10, inherits the selected joint recipe
and passes the seed-17 development screen on both encoders. Independent
audits verify its training certificates and development calculations.
Its composition-gain intervals still include zero, and the fixed three-seed
screen and matched controls are incomplete. This draft retains the failed
pilots and distinguishes this intermediate result from the practical
benchmark evidence needed for a performance or mechanism claim.

\section{Related work}
\paragraph{Non-regression in model updates.}
Yan et al.~\cite{pct} study negative flips, where an updated classifier
fails on examples its predecessor answered correctly, and use focal
distillation to emphasize those examples. Angioni et al.~\cite{rcat}
extend the concern to adversarial robustness and analyze learning with
non-regression constraints. Protecting baseline-correct decisions,
constrained retention, and the distinction between average gains and
individual regressions are therefore prior ideas. Our specialization
uses affine pairwise constraints over complete image--text galleries in
both directions, coupled to relation supervision and a finite numerical
audit. The earlier consistency results are not invoked as a generalization
theorem for our gallery-dependent retrieval procedure.

\paragraph{Frozen-feature alignment.}
LABCLIP learns a linear alignment of frozen image--text embeddings with
hard-negative contrastive supervision~\cite{labclip}. Its released scoring
implementation normalizes the transformed text. This is the closest
practical pooled-feature comparator. A learned linear map alone is not our
distinction; explicit complete-gallery ownership constraints and the joint
source/support objective are. Published results use different pretrained
checkpoints and training regimes from ours, so their numerical scores are
not substituted for matched measurements. BiCLIP also studies bilinear
adaptation, in few-shot classification~\cite{biclip}. These precedents
motivate a restrained architectural claim.

\paragraph{Composition and retained capabilities.}
CLIC combines compositional training examples with positive and negative
caption objectives and reports improvements on SugarCrepe++ and
retrieval~\cite{clic}. FSC-CLIP uses local hard-negative supervision and
selective calibration to improve composition while retaining multimodal
capabilities~\cite{fsc}. Thus retention is an established motivation,
not a new task introduced here. Our proposed comparison is between soft
or unconstrained adaptation and a checked finite-gallery constraint
mechanism. It requires equal-data controls; a favorable comparison to an
undertrained adaptation would not establish the mechanism's value.

\paragraph{Ranking and local alignment.}
RankCLIP uses listwise ranking consistency in pretraining~\cite{rankclip}.
Miranda et al. learn local alignments from frozen patch and token
representations~\cite{inference}. Our score instead uses cached pooled
features and protects specific source-ownership decisions. The respective
data regimes, representation access, and computational budgets differ.
We do not claim that pooled features are sufficient for all compositional
phenomena or that a ranking objective is itself new.

\paragraph{Recovering a frozen scoring endpoint.}
The author project page for CoSE describes a frozen channel and an
additional learned channel, with exact recovery of the frozen score at a
mixture endpoint~\cite{cose}. That is another reason not to claim novelty
for additive adaptation or baseline recovery alone. It differs from
protecting selected ownership decisions at a nonzero fitted state. This
citation is scoped to the author page: a paper and code were not linked at
the access date, and no venue or performance comparison is independently
established here.

\paragraph{Compatibility with stored gallery features.}
Backward-compatible training (BCT) learns updated representations that
can be compared with previously stored features, avoiding gallery
re-encoding during model updates~\cite{bct}. Reuse of an old gallery is
therefore an established deployment objective. Our query-only expression
is elementary algebra for the chosen bilinear score, with exact-real
equivalence in both retrieval directions. It neither introduces
backward compatibility in general nor establishes a new approximate
search guarantee.

\paragraph{Evaluation semantics.}
The e-ViL work introduces e-SNLI-VE and an explanation benchmark~\cite{evil}.
We use the relation annotations and associated image/source-caption
ownership to construct retrieval and relation-ranking evaluations.
Our R@1 measurements are therefore project-defined retrieval evaluations
on this data, not the original e-ViL explanation score. SugarCrepe++ tests
two semantically equivalent, lexically different positive captions against
a hard negative~\cite{sugarcrepepp}. Its role is deliberately distinct from
our source/support development proxy.

\section{Problem and method}
\subsection{Frozen embeddings and relation sets}
Let $v_i,t_j\in\R^d$ be the declared normalized frozen image and text
embeddings. Each source caption $j$ has one designated owning image $o(j)$.
The fitting gallery contains $n$ images and $M$ source captions, with
source-caption set $P_i$ for image $i$. In the present data, $|P_i|=5$.
Let $E_i$ contain supported descriptions and $N_i$ annotated contradictions.
The eligible composition owners are
$\mathcal I=\{i:|P_i|\,|E_i|\,|N_i|>0\}$.

We remove a contradiction from training supervision if its normalized
text equals any source or supported text for the same owner. Normalization
uses Unicode normalization, case folding, and word tokens with digits
retained. We do not convert neutral annotations into negatives. Such text
may contribute embeddings to the training-only PCA geometry. These rules
do not verify the correctness of every retained annotation.
Source-caption ownership is retained even when an identical string appears
under multiple images; its ambiguity is disclosed in Section~\ref{sec:data}.

\subsection{One centered residual score}
Training-only means $\mu_v,\mu_t$ and fixed orthonormal PCA bases
$U,V\in\R^{d\times r}$ define
\begin{align}
 x(v)&=U^\top(v-\mu_v),&y(t)&=V^\top(t-\mu_t),\\
 s_A(v,t)&=s_0(v,t)+x(v)^\top A y(t),&s_0(v,t)&=v^\top t.
 \label{eq:score}
\end{align}
Only $A\in\R^{r\times r}$ is learned, under $\norm{A}_F\le R$.
The $r=128$ model has 16,384 optimized parameters; the stored means and
bases are additional model state. All tasks and retrieval directions use
Equation~\eqref{eq:score}, with no inference-time labels, owner lookup, or
task-specific routing. The zero matrix gives the frozen baseline.
In exact arithmetic, the same score can also be implemented by transforming
only the query while retaining original frozen gallery vectors in either
direction (Appendix~\ref{app:query}). This is an elementary property of the
bilinear form, not a distinct model. Appendix~\ref{app:query} also reports
a limited arithmetic and warm dense-scoring diagnostic on training queries.

Centering removes an unconditional residual offset in the corresponding
finite pools: for a fixed text, averaging the residual over training images
gives zero, and the corresponding identity holds over the full PCA text
pool for a fixed image. The second pool includes neutral text and is not
the source-only retrieval gallery. Neither identity is a population
assumption. The model is an affine-score adaptation, not a claim to preserve
cosine geometry or a text-paraphrase representation.

\subsection{Full-gallery source retrieval}
Choose $p_i\in P_i$ as the best frozen-scoring owned caption, with ascending
manifest index resolving ties. It stays fixed during fitting. With frozen
positive logit scale $\kappa$, define
\begin{align}
 L_I(A)&=\frac1n\sum_i\log\left(1+\sum_{j:o(j)\ne i}
  \exp\{\kappa[s_A(v_i,t_j)-s_A(v_i,t_{p_i})]\}\right),\label{eq:i2t}\\
 L_T(A)&=\frac1M\sum_{j=1}^M\log\left(1+\sum_{k\ne o(j)}
  \exp\{\kappa[s_A(v_k,t_j)-s_A(v_{o(j)},t_j)]\}\right),\label{eq:t2i}\\
 L_{\mathrm{ret}}(A)&=\tfrac12[L_I(A)+L_T(A)].
\end{align}
Other source captions owned by image $i$ are excluded from its I2T
denominator, avoiding within-owner false negatives. Every source caption
is a T2I query. Query minibatches do not reduce either candidate gallery.
Holding the I2T anchor fixed makes each loss a log-sum-exp of affine scores.

\subsection{Joint relation ranking}
For $p\in P_i,e\in E_i,n\in N_i$, use
\begin{equation}
 \ell_{i,p,e,n}(A)=\log\left[
 1+\exp\!\left(\frac{s_A(v_i,t_n)-s_A(v_i,t_p)+\delta}{\tau}\right)
  +\exp\!\left(\frac{s_A(v_i,t_n)-s_A(v_i,t_e)+\delta}{\tau}\right)\right].
 \label{eq:joint}
\end{equation}
Average first over all triples within each image, then equally over
$\mathcal I$, to obtain $L_{\mathrm{comp}}$. This targets simultaneous
source and support correctness. It does not minimize a source--support
embedding distance or assert that the descriptions are equivalent.
All owners contribute source retrieval even if a relation set is missing.

The regularized objective is
\begin{align}
 F(A)&=L_{\mathrm{ret}}(A)+\lambda L_{\mathrm{comp}}(A)
           +\tfrac\eta2\norm{A}_F^2,\label{eq:objective}\\
 \lambda&=w\frac{\norm{\nabla L_{\mathrm{ret}}(0)}_F}
                    {\max(\norm{\nabla L_{\mathrm{comp}}(0)}_F,10^{-12})}.
 \label{eq:lambda}
\end{align}
The ratio is computed once from fitting data. It is not recomputed during
optimization or selected using official development. The expanded recipe
inherits $w=.25$; each encoder's actual multiplier depends on its training
gradients.

\subsection{Full-gallery retention constraints}
Let $Q_I$ and $Q_T$ be the image and source-caption training queries that
the frozen score answers correctly. For each $i\in Q_I$ and unowned caption
$j$, constrain its anchor margin:
\begin{equation}
 s_A(v_i,t_{p_i})-s_A(v_i,t_j)
 \ge\gamma[s_0(v_i,t_{p_i})-s_0(v_i,t_j)].\label{eq:ci}
\end{equation}
For each $j\in Q_T$ and $k\ne o(j)$, similarly require
\begin{equation}
 s_A(v_{o(j)},t_j)-s_A(v_k,t_j)
 \ge\gamma[s_0(v_{o(j)},t_j)-s_0(v_k,t_j)].\label{eq:ct}
\end{equation}
Here $0<\gamma\le1$; all experiments use $\gamma=.5$.
Write a frozen margin and its residual feature difference as $(m_a,D_a)$.
Then the feasible set is
\begin{equation}
 \mathcal C=\{A:\norm{A}_F\le R,\quad
                 \ip{A}{D_a}\ge-(1-\gamma)m_a\ \text{for all protected pairs }a\}.
 \label{eq:feasible}
\end{equation}
It includes zero and is a compact convex set. Both directions compare
against every irrelevant training competitor; no fixed shortlist defines
protection.

\begin{proposition}[Finite-gallery top-1 retention]\label{prop:retention}
For fixed galleries, ownership and features, use exact arithmetic and the
same ascending-index tie rule before and after adaptation. Choose each
anchor by frozen best-owned score with that rule. Every $A\in\mathcal C$
preserves correctness of every query in $Q_I$ and $Q_T$. Neither direction's
training R@1 decreases.
\end{proposition}
The proof, including frozen ties, appears in Appendix~\ref{app:proof}.
This proposition protects correctness, not every ranking position or the
identity of a winning owned caption. New owners, new gallery competitors,
other recall cutoffs, and other data distributions are outside its scope.
The existence of zero in $\mathcal C$ does not establish a useful nonzero
solution. Our experimental gates explicitly reject a zero update.

\subsection{Optimization and numerical verification}
We initialize $A=0$, use seeded image-query minibatches, and combine the
objective gradient with an active-constraint squared-violation penalty.
Each update clips the gradient and projects onto the Frobenius ball. At
epoch end, full-gallery scanning identifies violations, directional passes
repair active inequalities, and an algebraically computed contraction
toward zero restores feasibility. The contraction is based on training
constraints, not a development amplitude search. This repair is not claimed
to be the closest-point projection.

An accepted constrained checkpoint must satisfy all inequalities within
$10^{-12}$ and lose zero protected queries under the canonical scorer.
The scorer uses fixed float64 pair reductions; exhaustive matrix screening
with conservative error bounds identifies candidates requiring canonical
rescoring around near ties. The checkpoint with the lowest training
objective among nonzero feasible states is selected; ties favor the earlier
epoch. Official development selects neither epoch nor residual amplitude.

The V10 implementation stores the full frozen similarity matrix in two
immutable, disk-backed orientations and streams query blocks. Every
candidate remains in the loss and constraint scan. This changes storage
and floating-point reduction order, not the mathematical gallery or
objective. Prefit comparisons to the dense implementation and independent
gradient checks passed. Numerical agreement does not imply bitwise
identity of entire fitted trajectories.

The objective is strongly convex for $\eta>0$, but our finite stochastic
optimizer is not proved to reach its unique constrained optimum.
Appendix~\ref{app:proof} gives the conservative reported bound
$\ip{\nabla F(A)}{A}+R\norm{\nabla F(A)}_F$. It concerns training objective
suboptimality, not generalization.

\section{Data and experimental protocol}\label{sec:data}
\subsection{Encoders and fitting pools}
We use two pinned, heterogeneous checkpoints: ViT-B/32 trained on LAION2B,
with 512-dimensional embeddings, and OpenAI RN50, with 1,024-dimensional
embeddings. The former is not the OpenAI ViT-B/32 used in several published
comparisons. Images use the pinned 224-pixel RGB preprocessing and texts
the pinned tokenizer/causal encoder. A declared float32 normalization pass
precedes float64 adaptation math. The expanded exports preserve the original
1,200 image and 22,192 text feature rows bitwise; metadata binds encoder
weights, preprocessing, normalization, and feature hashes.

Relation labels come from official e-SNLI-VE training annotations. V8/V9
fit 1,200 owners. V10 adds 4,800 owners selected by a fixed SHA-256 ranking
of eligible official-training IDs, excluding reserved IDs. Training labels
had been materialized before the sample was frozen, but selection does not
use labels, features, or measured outcomes. Table~\ref{tab:data} gives the
raw annotation inventory. These row counts are not independent statistical
sample sizes. Neutral rows enter geometry only.

\begin{table}[t]
\centering
\caption{Fitting-data inventory. Counts are raw image-associated text rows;
loss eligibility follows the stated conflict rule. V10's two seed-17 fits
are complete; replication remains ongoing.}\label{tab:data}
\begin{tabular}{lrr}
\toprule
Quantity & V8/V9 & V10 frozen pool\\
\midrule
Image owners & 1,200 & 6,000\\
Source captions & 6,000 & 30,000\\
Supported descriptions & 5,319 & 26,374\\
Contradictions & 5,846 & 29,007\\
Neutral descriptions & 5,027 & 25,135\\
Total text rows & 22,192 & 110,516\\
Distinct exact strings & 21,744 & 104,963\\
Source strings shared across owners & 0 & 22\\
\bottomrule
\end{tabular}
\end{table}

The 22 shared source strings in V10 retain their original ownership labels;
we do not silently relabel equivalence or remove them. They can create
false negatives and tie ambiguity. Exact normalized contradiction conflicts
are filtered only as declared. We add no annotation-quality filter claiming
that all remaining supports or contradictions are correct.

Image audits decoded all 7,500 selected training and fresh-reserve images.
No exact JPEG or decoded RGB duplicates were found within those selected
owners. Exact JPEG hashes were checked against 8,377 earlier reserved
owners; recompressed or perceptually similar images across those earlier
reserves were not exhaustively excluded. Consequently, owner disjointness
and these exact checks do not establish absence of all content overlap or
pretraining contamination.

\subsection{Development endpoints and selection}
Official development retrieval uses a fixed gallery of 900 images and
4,500 owned source captions. I2T R@1 is correct if the top caption is owned
by the image; T2I R@1 is correct if the top image owns the caption.
Composition development uses 100 image clusters. The original joint metric
is the within-image mean of
\begin{equation}
 \indicator\{\min(s(v_i,t_p),s(v_i,t_e))>s(v_i,t_n)\},
 \quad p\in P_i,e\in E_i,n\in N_i,
\end{equation}
then averaged equally over eligible images. A source-pair companion replaces
the supported description by a distinct source caption. Original joint
evaluation retains all raw annotated labels rather than applying the
training eligibility filter; the source-pair metric uses the fixed
normalized-conflict and distinct-caption rules. Both use strict positive
comparisons and never treat neutral labels as contradictions.

V9 chose its configuration using one owner-disjoint 960/240 split of the
original training pool, shared by both encoders. Only 960 owners contributed
fitting and PCA; validation queries came from the 240 held-out owners.
The validation retrieval gallery nevertheless contained all 1,200 images
and 6,000 source captions. The finite grid used $R\in\{.1,.3,1\}$ and
$w\in\{.25,1,4\}$ at rank 128. Eligibility required a strictly positive
original-joint mean gain and nonnegative changes in the other three
endpoints on both encoders. The common configuration maximized the worst
encoder's original-joint gain, followed by declared tie rules. Only
$(R,w)=(1,.25)$ was eligible. This adaptive inner selection is not
confirmatory evidence.

V10 inherits that one configuration with no new inner grid. Means, bases,
and the fixed gradient calibration are recomputed from its training pool.
Table~\ref{tab:recipe} distinguishes the earlier and expanded recipes.
V8/V9 are successive research attempts, not controlled single-factor
ablations of each other.

\begin{table}[t]
\centering\small
\caption{Locked optimization recipes. Update counts are complete fit
budgets, not necessarily the selected epoch's update count.}\label{tab:recipe}
\begin{tabular}{lrrr}
\toprule
Setting & V8 & V9 & V10\\
\midrule
Owners & 1,200 & 1,200 & 6,000\\
Rank $r$ & 64 & 128 & 128\\
Radius $R$ & .1 & 1 & 1\\
Epoch budget & 24 & 32 & 32\\
Batch size & 64 & 64 & 64\\
Optimizer update budget & 456 & 608 & 3,008\\
Base learning rate & .01 & .03 & .03\\
Source retrieval loss & absent & present & present\\
Composition positives & support & source/support & source/support\\
Weight calibration $w$ & not used & .25 & .25\\
\bottomrule
\end{tabular}
\end{table}
The learning rate decreases as the inverse square root of the epoch.
Shared settings are $\gamma=.5$, $\tau=.02$, $\delta=.005$, ridge $.01$,
active penalty weight 100, gradient clipping at 10, at most 12 repair
rounds with three directional passes, and tolerance $10^{-12}$.
The V10 fit uses three CPU threads and query blocks of 64.
Moving from 608 to 3,008 updates changes data and optimization together,
along with gallery hardness and PCA geometry. It is an expanded
data-and-compute recipe, not an isolated data-size or sample-efficiency
experiment. Feature extraction and adapter fitting costs must be reported
separately.

\subsection{Uncertainty and unchanged continuation gates}
Let $D_{is}$ be the paired trained-minus-frozen outcome for image cluster
$i$ and seed $s$; for T2I it averages the owned caption-query outcomes within
the image. The fixed three-seed summary is the equal mean of per-query
outcomes for seeds 17, 29, and 43. It is not a score or feature ensemble.
We resample paired image clusters with 100,000 draws, bootstrap seed
20261007, nominal familywise level $.05$, and family size 80. The two-sided
tail probability is $.05/(2\cdot80)=.0003125$.

SugarCrepe++ has a different aggregation unit from the image-balanced
development proxy: the target score averages declared caption-triplet
outcomes over 4,757 items associated with 1,542 images in the fixed
benchmark manifest. When image clusters are resampled, both their summed outcome
differences and their numbers of triplets are resampled. The difference
is their ratio, with a random total-item denominator. It is not the mean
of equally weighted image-level triplet accuracies. Image clusters still
preserve dependence between captions belonging to the same image.

Intervals condition on fixed selected states and a fixed full gallery.
They resample query-cluster outcomes without reranking a randomly sampled
gallery. They do not measure variation over new training seeds, new
galleries, or adaptive model selection. Nominal multiplicity adjustment
does not correct the whole historical search, including earlier reuse of
the practical benchmark data. The later locked practical comparison
program contains 60 effects: 30 core, 20 retrieval-only-control, and 10
expanded LABCLIP-style contrasts. The fresh same-source program contains
24 effects in a separate declared family. Each uses the inherited
family-80 adjustment; this is not a pooled simultaneous 95-percent
statement over all 84 effects.

For each encoder, the seed-17 development screen requires a nonzero update,
positive original-joint mean gain, nonnegative source-pair mean gain,
nonnegative mean changes in I2T and T2I, and both adjusted retrieval lower
bounds strictly greater than $-.01$. Both encoders must pass before
replication. Seeds 29 and 43 then use the unchanged recipe; the same
development contract applies to the fixed three-seed mean. Individual
replication seeds are not selected or required to pass separately.
A failed stage stops the recipe without an official-development checkpoint
or amplitude search.

Only after the three-seed development gate passes does the practical
contract permit a locked final evaluation. It requires both encoders,
nonzero trained states for all three seeds, e-ViL-derived test retrieval
over the full 1,000-image/5,000-caption gallery with I2T and T2I adjusted
lower bounds strictly above $-1$ percentage point, and SugarCrepe++
both-positive correctness gain with adjusted lower bound strictly above
zero. COCO full-gallery retrieval is secondary.
\mbox{Equality fails a strict bound.} No V8/V9 final benchmark result is
introduced here.

\subsection{Baseline and matched controls}
The V9 LABCLIP-style comparator uses a full-rank bias-free $W$, initialized
at identity, and the universal score $v^\top\mathrm{normalize}(Wt)$.
It uses source positives and one valid owner-associated contradiction;
supported labels are consulted for exact conflict exclusion but are not
direct positive targets. Owners with no valid contradiction are omitted
and recorded, while our source retrieval uses all owners. These are
important supervision differences.

The published-style control uses batch 1,024, Adam learning rate $.001$,
10 epochs, and a learned log scale initialized to zero. A separately named
small-data adaptation uses owner-unique batches of 128, the fixed native
logit scale, and five source-caption passes per epoch over 32 epochs.
Its candidate epochs are $1,2,4,8,16,32$. Identity-map normalization is
evaluated separately from the original frozen scores. We do not use an
undertrained published-style control as a performance ceiling.

The V10 protocol binds a no-retention control with the same data, geometry,
objective, calibrated weight, initialization, paired query orders, norm
radius, optimizer budget, and objective-based checkpoint rule. It removes
the active retention penalty and feasibility repair; it does not require
checkpoint feasibility. Training losses and protected-query regressions
remain diagnostics. This contrast identifies the entire enforcement
procedure, not hard constraints separately from their optimizer.
Control fits are scheduled only after both joint seed-17 development
pilots pass. They are not completed results in this manuscript.
Any eventual paired effect is conditional on this development-selected
recipe, not an unconditional causal benefit across all model recipes.

\section{Completed results}
\subsection{Training certificates and development transfer}
V8 and V9 both preserve all frozen-correct training queries for both
encoders. Table~\ref{tab:certificate} reports the protected counts shared
by these fits. V9's selected states were independently checked over
20,260,702 inequalities and 14,400,000 full-gallery image--text pairs across
the two encoders. The recorded numerical checks use tolerance $10^{-12}$
and independently require zero lost canonical correct queries. Their scope
is the fitting gallery.

\begin{table}[t]
\centering
\caption{Training retention for the 1,200-owner fitting gallery.
Both V8 and V9 have zero lost protected top-1 queries.}\label{tab:certificate}
\begin{tabular}{lrr}
\toprule
Quantity & ViT-B/32 LAION2B & RN50 OpenAI\\
\midrule
Protected I2T queries & 998 & 922\\
Protected T2I queries & 3,914 & 3,384\\
Checked inequalities per state & 10,675,896 & 9,584,806\\
V8 selected coefficient norm & .099934 & .099917\\
V9 selected coefficient norm & .999191 & .999167\\
V8 selected epoch & 24 & 23\\
V9 selected epoch & 32 & 32\\
\bottomrule
\end{tabular}
\end{table}

\begin{samepage}
Table~\ref{tab:means} shows the completed official-development point
estimates. Frozen R@1 was 85.000/68.200 percent (I2T/T2I) for ViT-B/32 and
81.222/58.422 percent for RN50. The frozen original-joint proxy was
81.955 and 78.648 percent, respectively.\par
\end{samepage}

V8's supported-only adaptation
has negligible composition changes and negative I2T means. V9 has positive
means for every endpoint, with its largest change in RN50 T2I.
V10 also has positive means for all four endpoints and passes the initial
development screen on both encoders. It does not yet supply a replicated
effect or a practical benchmark result.

\begin{table}[t]
\centering\small
\caption{Seed-17 official-development paired changes, in percentage
points. Original joint and source-pair joint are development proxies,
not SugarCrepe++ scores. V8/V9 fail their development screens; V10 passes
the seed-17 screen only.}
\label{tab:means}
\begin{tabular}{llrrrr}
\toprule
Recipe & Encoder & I2T R@1 & T2I R@1 & Original joint & Source-pair joint\\
\midrule
V8 & ViT-B/32 & $-.3333$ & $-.0889$ & $+.0319$ & $-.2062$\\
V8 & RN50 & $-.4444$ & $+.1333$ & $-.0262$ & $-.1698$\\
V9 & ViT-B/32 & $+.6667$ & $+.9333$ & $+.8079$ & $+.1047$\\
V9 & RN50 & $+1.1111$ & $+3.2444$ & $+.8553$ & $+.0167$\\
V10 & ViT-B/32 & $+1.0000$ & $+.9111$ & $+.6701$ & $+.0185$\\
V10 & RN50 & $+1.4444$ & $+4.0889$ & $+1.0788$ & $+.0190$\\
\bottomrule
\end{tabular}
\end{table}

The favorable V9 means do not satisfy the retention requirement. Its I2T
adjusted lower bounds are $-1.528$ and $-1.778$ percentage points,
below the $-1$ threshold (Table~\ref{tab:ci}). The two composition-gain
intervals also include zero. The development screen requires positive
composition means, whereas the final practical composition gate would
require a strictly positive adjusted lower bound. Neither inference can
be replaced by inspecting favorable point estimates.

\begin{table}[t]
\centering\small
\caption{Nominal family-80 adjusted paired bootstrap intervals in
percentage points. Intervals condition on the selected seed-17 state and
fixed gallery.}\label{tab:ci}
\begin{tabular}{llrrr}
\toprule
Recipe & Encoder & I2T R@1 & T2I R@1 & Original joint\\
\midrule
V8 & ViT-B/32 & $[-1.444,.667]$ & $[-.489,.311]$ & $[-.928,.819]$\\
V8 & RN50 & $[-1.667,.556]$ & $[-.356,.644]$ & $[-.899,.509]$\\
V9 & ViT-B/32 & $[-1.528,2.889]$ & $[-.333,2.200]$ & $[-.597,2.293]$\\
V9 & RN50 & $[-1.778,4.111]$ & $[1.672,4.844]$ & $[-.980,2.607]$\\
V10 & ViT-B/32 & $[-.778,2.889]$ & $[-.311,2.133]$ & $[-.581,1.854]$\\
V10 & RN50 & $[-.889,3.889]$ & $[2.689,5.528]$ & $[-.426,2.506]$\\
\bottomrule
\end{tabular}
\end{table}

The distinction between positive mean and retention evidence is visible
in the paired I2T counts. V9 corrects 20 cases and regresses on 14 for
ViT-B/32, a net six out of 900; RN50 corrects 35 and regresses on 25, a
net ten. The uncertainty depends on this turnover as well as its net
effect. Increasing training seeds does not automatically remove it:
per-image seed outcomes can be correlated. No replication occurred after
these failed screens. We do not use the changed cases for case-specific
development constraints.

\subsection{What the normalized baseline establishes}
No common candidate epoch of the adapted LABCLIP-style baseline passed the
V9 inner gate. All six tested epochs decreased inner T2I R@1 for each
encoder. The family therefore stopped before full fitting and official
development. This result is evidence about the declared recipe and one
inner split, not general inferiority of LABCLIP.

\begin{samepage}
Different normalization,
capacity, positive targets, owner eligibility, and update counts prevent
calling it a single-factor ablation of the constrained method. In
particular, V9's baseline outcome does not establish its behavior on the
expanded V10 pool.\par
\end{samepage}

\subsection{Expanded V10 seed-17 results}
Both 6,000-owner fits select epoch 32 by their training-only objective,
after 3,008 updates. Independent full-fit audits verify 64 epoch checkpoint
files (epochs 1--32 for both encoders) and the recorded minimum selection,
and recompute the two selected
states' PCA geometry, objectives, canonical winners and constraints.
They check 376,071,311 inequalities and 360,000,000 selected-state
image--text pairs across the encoders. No protected training query is
lost (Table~\ref{tab:v10certificate}). Objective eligibility uses 5,979
relation owners and 746,320 triples; every owner contributes retrieval.
This is a completed numerical certificate for these two states, not a
claim that all epoch objectives were independently replayed or that the
finite optimizer reached the unique optimum.

\begin{table}[t]\centering\small
\caption{V10 seed-17 training states on the 6,000-image/30,000-caption
gallery. Canonical correct-query counts and retention are independently
audited with constraint tolerance $10^{-12}$. The multiplier and relaxed
bound are recorded diagnostics; their gradient calculations were not
independently replayed.}\label{tab:v10certificate}
\begin{tabular}{lrr}
\toprule
Quantity & ViT-B/32 LAION2B & RN50 OpenAI\\
\midrule
Frozen correct I2T / T2I & 3,948 / 14,244 & 3,565 / 10,880\\
Adapted correct I2T / T2I & 4,295 / 15,435 & 3,938 / 12,813\\
Lost protected I2T / T2I & 0 / 0 & 0 / 0\\
Checked inequalities & 203,870,016 & 172,201,295\\
Coefficient Frobenius norm & .986485 & .983216\\
Fixed multiplier $\lambda$ & .869346 & 1.656168\\
Selected training objective & 2.258950 & 3.090089\\
Ball-relaxed bound & .086495 & .224596\\
Selected epoch / updates & 32 / 3,008 & 32 / 3,008\\
\bottomrule
\end{tabular}
\end{table}

The independently reconstructed seed-17 development effects appear in
Tables~\ref{tab:means}--\ref{tab:ci} and Figure~\ref{fig:v10development}.
Both retrieval means increase, with adjusted I2T lower bounds of $-.778$
and $-.889$ percentage points, respectively. These exceed the $-1$-point
retention threshold by only $.222$ and $.111$ points. ViT-B/32 corrects
16 I2T queries and regresses on 7; RN50 corrects 26 and regresses on 13.
For T2I the corresponding improved/regressed counts are 116/75 and 244/60.
Thus the development result is compatible with individual regressions;
the zero-regression certificate applies only to training.

The original-joint gains of $+.6701$ and $+1.0788$ points and the small
positive source-pair changes satisfy the mean-based composition screen.
Both original-joint and source-pair gain intervals nevertheless include
zero. The development screen must not be reported as statistically
established composition improvement, and it is not SugarCrepe++.
The independent development audit checks 16,200,000 canonical pairs,
83,880 frozen/adapted composition triples, and 800,000 bootstrap draws
across both encoders. Retrieval means and interpolated interval endpoints
are also checked with exact integer/rational accounting. Replication
with the fixed seeds 29 and 43 is underway; its aggregate gate is incomplete.

\begin{figure}[!htbp]\centering
\includegraphics[width=\linewidth]{v10_development_effects.png}
\caption{Audited V10 seed-17 development effects. The initial screens pass,
but the composition intervals cross zero. Intervals condition on these
selected states and fixed galleries; no practical benchmark result or
three-seed conclusion is shown.}\label{fig:v10development}
\end{figure}

\clearpage
\section{Ongoing expanded experiment and missing comparisons}
\label{sec:ongoing}
\paragraph{Frozen V10 experiment.}
Feature exports, the 6,000-owner manifest, inherited recipe, numerical
implementation, and prefit audit remain locked. The two seed-17 fits and
development audits are complete and pass their initial screens. The next
scientific decision is the unchanged two-encoder, fixed-three-seed
development gate. No complete replication aggregate, matched-control
effect, practical benchmark result, or fresh-reserve outcome is reported.
The inherited recipe, epoch budget and checkpoint rule remain unchanged.

\paragraph{Fresh same-source confirmation.}
A separately salted fixed lock reserves 1,500 unused official-training
owners, forming a future complete 1,500-image/7,500-caption gallery.
No fitting or selection uses these owners. The contract forbids their
feature encoding until all 12 core final states (joint and no-retention,
two encoders, three seeds) and the evaluation sources are frozen. A
separate release rule additionally requires all 24 planned states,
including the retrieval-only and LABCLIP-style controls, before fresh
outcome scoring. Only the two original families enter fresh contrasts.
Its contrasts are joint minus frozen, no-retention minus frozen, and joint
minus no-retention, each on I2T, T2I, original joint and source-pair joint
for both encoders. This gives 24 declared endpoint contrasts in the
separate fresh family-80 calculation. All are to be reported; they are not
pooled with the historical-benchmark family's 60 contrasts.

The supplementary joint-versus-frozen criterion requires both retrieval
lower bounds $>-.01$, original-joint gain lower bound $>0$, and nonnegative
source-pair mean change, for both encoders. It remains a same-source,
owner-held-out check. It cannot substitute for SugarCrepe++ or the required
full-gallery practical retrieval test, cannot prove a new-domain claim,
and does not retroactively make historical benchmark reuse confirmatory.

\paragraph{Evidence needed for mechanism and performance claims.}
The locked joint/no-retention contrast is necessary to attribute a change
to the enforcement procedure. If both arms behave similarly outside
training, the certificate remains a property of the constrained fit but
does not establish an empirical advantage. A constraint benefit alone
does not show that supported annotations matter.

A separate explanatory retrieval-only contract was frozen before
V10 held-out outcomes. It uses the same 6,000 owners, complete source
gallery, retention procedure, rank, radius, 32 epochs, and fixed seeds
17, 29, and 43, with $\lambda=0$. It selects its own minimum nonzero
feasible training-objective checkpoint. This arm is explanatory, never a
candidate for model selection, and runs only if the joint three-seed
development gate passes and all six no-retention fits are complete. It
adds joint-minus-retrieval-only and retrieval-only-minus-frozen contrasts
on five practical endpoints and two encoders, giving 20 effects. It is
not added to the two-family fresh-1,500
confirmation contract. Its results are not available here.
A composition-only arm at matched capacity is still needed to isolate
the full-gallery retrieval objective; V8 is not that matched arm.

To test the value of support labels, a count-matched source-positive
control should retain per-owner term counts and negative exposure while
replacing supported slots with a frozen seeded draw of source captions.
This matches weighted positive exposure, but not the number of distinct
descriptions, and must be interpreted accordingly. This support control
and the composition-only control remain proposed evidence requirements,
not locked or completed V10 arms.

\newpage
\paragraph{Expanded normalized alignment comparator.}
A separately locked explanatory LABCLIP-style arm uses the same 6,000-owner
feature pool, pure normalized full-rank map, fixed native logit scale,
owner-unique batch size 128, 32 epochs, five source passes per epoch, and
seeds 17, 29, and 43. Adam uses learning rate $.001$, betas $(.9,.999)$,
epsilon $10^{-8}$ and zero weight decay. Owners without a valid contradiction
are omitted and must be counted. The annotated owner contradictions are a
declared supervision adaptation, not the published word-shuffle recipe;
supported labels enter conflict exclusion only. Its possible 7,520 updates
are not the joint method's 3,008, and the parameter counts also differ.

Each epoch is evaluated on the same seeded, owner-unique HNB audit batches
covering five source rows per eligible owner, with one fixed contradiction
per row. The objective is a row-weighted fixed-batch diagnostic, not
full-gallery or all-negative contrastive loss. Selection takes its earliest
minimum among functionally nonzero epochs; movement is measured against
the identity-normalized map on fixed training witnesses. Identity-map
versus original-frozen numerical parity is reported separately. No held-out outcome
selects the epoch, normalization or scale.

The arm runs only after the joint three-seed development gate passes,
cannot replace the primary candidate, and never enters fresh-1,500
contrasts. All six states must be frozen before practical benchmark scores
are released. Its ten paired joint-minus-LABCLIP-style contrasts cover the
five practical endpoints and two encoders; baseline absolute metrics are
descriptive. No expanded baseline result or superiority claim is available.

A sample-efficiency claim additionally requires an update-matched
learning-curve design; 1,200 versus 6,000 owners at equal epochs does not
isolate data. A broad comparison to CLIC/FSC would require matched
pretrained checkpoints and transparent data/compute accounting, or a
clearly delimited claim about pooled-feature methods. Published
cross-checkpoint numbers are contextual evidence only.

\section{Limitations and research integrity}
\paragraph{No established practical solution.}
V8 and V9 fail the continuation rule. V10 passes its initial development
screen, but its composition-gain intervals cross zero and its replication
aggregate remains incomplete. A passed development screen does not
establish the practical benchmark goal. Nothing in the theorem guarantees
a useful nonzero feasible point,
an improved composition score, or transfer of retained training decisions.

\paragraph{Restricted geometry and supervision.}
Pooled embeddings may omit information needed for localized composition.
Rank-restricted residuals cannot create information absent from those
features. PCA favors high-variance directions rather than necessarily
composition-relevant directions. Support annotations do not identify
paraphrases. Ownership labels may mark a semantically compatible caption
as irrelevant to another image. The duplicate-source issue and remaining
annotation noise are preserved limitations, not removed by exact conflict
filtering.

\paragraph{Finite numerical scope.}
The mathematical proposition uses exact arithmetic. Our audit uses a
declared tolerance plus direct deterministic rankings. Strong convexity
does not imply exact convergence of the implemented repair-based optimizer.
Its feasible set protects baseline-correct queries only; the theorem is
neither a zero-shot capability guarantee nor protection against a changed
gallery. Disk-backed exhaustive training scans have quadratic gallery-storage
cost. The query-only deployment identity avoids re-encoding the gallery,
but the measured subset does not certify arbitrary floating-point or
approximate retrieval implementations.

\paragraph{Adaptive evidence and dataset overlap.}
The current program follows earlier model attempts and historical benchmark
reuse. The inner split, official development set, and final benchmark
roles are now explicit, but the family-80 bootstrap adjustment does not
cover all previous adaptive decisions. We describe reused benchmark
results as exploratory even if a later locked recipe meets its numerical
gate. The fresh same-source reserve supplies a different, supplementary
check after freezing all choices; exact image checks do not establish
absence of near duplicates or overlap with encoder pretraining.

\paragraph{Incomplete comparison set.}
The matched enforcement-removal result, objective removals, support-count
control, expanded normalized baseline, and clean final evidence are not
available. Without them, attributing improvements to constraints or
annotations would be premature. No original reviewer text or recovered
manuscript is quoted here, and no claim is made that these experiments have
resolved particular review comments. This draft is a new synthesis of
verified code, protocols, aggregate results, and primary literature.

\section{Conclusion}
Full-gallery constraints make a limited but testable promise: they retain
every frozen-correct training top-1 decision under a shared score. The
completed experiments verify that promise and show why it is insufficient
for practical transfer. The smaller joint pilot improves development means
yet fails the retention bound. The expanded recipe passes the initial
two-encoder development screen, while its composition intervals still
include zero. Replication and locked controls must establish whether the
formulation offers an advantage beyond this intermediate result. The
current contribution is the constrained formulation, its finite-gallery
analysis, and audited training and development evidence; practical
compositional improvement with retained retrieval remains unestablished.

\appendix
\section{Proofs and optimization scope}\label{app:proof}
\paragraph{Proof of Proposition~\ref{prop:retention}.}
For an I2T constraint,
$D_a=x(v_i)[y(t_{p_i})-y(t_j)]^\top$; for T2I,
$D_a=[x(v_{o(j)})-x(v_k)]y(t_j)^\top$.
The adapted margin is $m_a+\ip{A}{D_a}\ge\gamma m_a$.
Frozen correctness implies $m_a\ge0$. A strictly positive frozen margin
therefore stays positive. If $m_a=0$, the fixed frozen-correct anchor
precedes the irrelevant competitor under the same index tie rule. The
adapted margin is nonnegative, so either it becomes positive or the
unchanged tie rule still favors that anchor. Applying this argument to
every irrelevant competitor preserves each protected query. Other
queries may improve or stay incorrect; none can reduce the already
preserved count. \hfill$\square$

\begin{proposition}[Convexity and a relaxed suboptimality bound]
For fixed features, anchors, $\lambda\ge0$, and $\eta>0$, $F$ is
$\eta$-strongly convex and has a unique minimizer $A^*$ on $\mathcal C$.
For feasible $A$ and $G=\nabla F(A)$,
\begin{equation}
 0\le F(A)-F(A^*)\le\ip{G}{A}+R\norm{G}_F.
\end{equation}
\end{proposition}
\begin{proof}
Every unregularized term is a log-sum-exp of affine functions of $A$.
Adding $\eta\norm{A}_F^2/2$ yields strong convexity. The feasible set is
nonempty, convex and compact, giving existence and uniqueness of the
minimizer. Convexity gives
$F(A)-F(A^*)\le\ip{G}{A-A^*}$. Relaxing $A^*\in\mathcal C$ to
$\norm{A^*}_F\le R$ yields the ball support function and the displayed
bound. No Slater condition or strict positivity of every frozen margin is
needed for this bound.
\end{proof}
The unique optimum can be zero; the result does not establish nonzero
feasibility or practical utility. V9's selected numerical bound values are
approximately .009253 (ViT-B/32) and .012375 (RN50); V10 seed-17 values are
.086495 and .224596. They are conservative
ball-relaxed bounds, not an exact constraint duality gap. Floating-point
gradients and tolerance limit their numerical interpretation.

For a new I2T anchor/competitor, Cauchy--Schwarz gives
\begin{equation}
 |[s_A(v,t_p)-s_A(v,t_n)]-[s_0(v,t_p)-s_0(v,t_n)]|
 \le R\norm{x(v)}_2\norm{y(t_p)-y(t_n)}_2.
\end{equation}
A frozen margin exceeding this quantity remains positive. The analogous
T2I bound exchanges the image difference and text coordinate. These
deterministic sufficient conditions do not bound population retention
without assumptions on unseen margins and features.

\subsection{Exact-arithmetic query-only scoring}\label{app:query}
Let $B=UAV^\top\in\R^{d\times d}$. Define
\begin{align}
 q_I(v)&=v+B^\top(v-\mu_v),
 &c_I(v)&=-(v-\mu_v)^\top B\mu_t,\\
 q_T(t)&=t+B(t-\mu_t),
 &c_T(t)&=-\mu_v^\top B(t-\mu_t).
\end{align}
Then the one score in Equation~\eqref{eq:score} satisfies
\begin{equation}
 s_A(v,t)=q_I(v)^\top t+c_I(v)
         =v^\top q_T(t)+c_T(t).
 \label{eq:query-scoring}
\end{equation}
To verify the first equality, expand
$(v-\mu_v)^\top B(t-\mu_t)$ into its term linear in $t$ and its
remaining query-only term. Expanding instead in $v$ gives the second
equality. This is elementary bilinear algebra.

Consequently, exact maximum-inner-product retrieval can use the original
frozen $d$-dimensional gallery vectors in both directions. For each query,
the corresponding offset is constant across candidates and does not affect
their order; adding it recovers the raw score. Both query transforms use
the same $A,U,V,\mu_v,\mu_t$ and implement one function, rather than
task-conditioned parameter sets. The factored transform requires
$O(dr+r^2)$ arithmetic per query, or $O(dr)$ with $UA$ precomputed.
These are algebraic operation counts, not latency or memory guarantees.

This observation gives compatibility with unchanged frozen gallery vectors
under exact inner-product search. It does not certify an approximate
nearest-neighbor index or any index-specific search procedure. Nor does
the identity guarantee bitwise equality after floating-point reassociation:
the numerical certificates in this study apply to the declared canonical
scorer and tie rule. A different implementation needs its own numerical
agreement and retrieval checks before inheriting a practical certificate.

\paragraph{Measured training-query diagnostic.}
For each audited V10 seed-17 state, we take 32 predetermined, evenly spaced
training queries in each direction and compare query-transformed dense
scores against the canonical scorer over the unchanged full gallery.
Across the two encoders this covers 2,304,000 pairs. All top-1 choices
agree, both with and without restoring the query-only offset; the largest
observed restored-score error is $7.22\times10^{-16}$. These are observations
on this fixed subset, not a worst-case floating-point bound or a replacement
for the training certificate.

\begin{table}[t]\centering\small
\caption{Warm CPU dense-scoring diagnostic, batch size 32. Ratios are
median query-transformed scoring time divided by frozen scoring time;
they exclude feature encoding, I/O, ANN, top-$k$ extraction and canonical
rescoring.}\label{tab:querytiming}
\begin{tabular}{llrr}
\toprule
Encoder & Direction & Gallery vectors & Time ratio\\
\midrule
ViT-B/32 & I2T & 30,000 & $1.021\times$\\
ViT-B/32 & T2I & 6,000 & $1.059\times$\\
RN50 & I2T & 30,000 & $1.012\times$\\
RN50 & T2I & 6,000 & $1.108\times$\\
\bottomrule
\end{tabular}
\end{table}
The timing experiment uses float64 dense scoring on the available x86-64
CPU with one BLAS thread, three warmups and 11 measured repeats, with no
co-scheduled fitting job. The end-to-end diagnostic includes the factored
query transform and dense scoring, after fixed model setup. Queries are
not renormalized. Table~\ref{tab:querytiming} reports batch-cost ratios,
not single-query latency, encoder inference costs, hardware-independent
speed claims, or measured approximate-index behavior.

\section{Reproducibility records and exact result provenance}
All numerical tables above were transcribed from immutable aggregate
result records, not reconstructed by opening benchmark cases. The
companion \path|practical_v9_claims_evidence_20261005.json| contains
the detailed V9 evidence map. The following abbreviated SHA-256 prefixes
identify central records; full hashes remain in the respective protocols.
\begin{center}\small
\begin{tabular}{p{.60\linewidth}l}
\toprule
Record & SHA-256 prefix\\
\midrule
V8 constrained protocol & \texttt{69ce998049bb}\\
V9 alignment protocol & \texttt{9e1bd80f901c}\\
V9 shared inner selection & \texttt{305e5277a8dd}\\
V9 independent full-fit audit & \texttt{c227361f3886}\\
V10 expanded training protocol & \texttt{ce04032395f1}\\
V10 fresh confirmation contract & \texttt{b1f426b28492}\\
V10 prefit audit & \texttt{32b4d54f7fc1}\\
V10 seed-17 independent full-fit audit & \texttt{2b2a075cdb57}\\
V10 seed-17 independent development audit & \texttt{f6daabe669a5}\\
V10 ViT-B/32 development result & \texttt{09dba4d91b66}\\
V10 RN50 development result & \texttt{696080db4157}\\
V10 retrieval-only contract & \texttt{ea3b5cb07328}\\
V10 replacement LABCLIP-style contract & \texttt{9772bec82d15}\\
V10 ViT-B/32 query-transform diagnostic & \texttt{0c3a1220045d}\\
V10 RN50 query-transform diagnostic & \texttt{e749a5ac0be4}\\
Expanded caption inventory & \texttt{d6de04798e5b}\\
\bottomrule
\end{tabular}
\end{center}
V8 aggregate effects come from
\path|results/practical_v8/development/{encoder}_seed17/result.json|;
V9 effects come from
\path|results/practical_v9/official_development/{encoder}_seed17/result.json|.
The V8 independent development audit reconstructed all 800,000 bootstrap
values for the two encoders. The V9 full-fit audit independently recomputed
selected-state objectives, canonical retrieval winners, and inequalities.
Audit success means artifact or numerical correctness, not empirical
success of the practical contract.

The V10 prefit audit checked 52 bound sources and verified unchanged
earlier sources, the locked sampling, annotation reconstruction, encoder
identity, old-row feature parity, full-gallery numerical equivalence,
and control implementation. Its test report has 29 passing tests and
16 passing subtests. A preserved profiling timer correction changes only
timing accounting, not numerical results. These records precede fitting;
they are distinct from the later model audit. The V10 full-fit audit
subsequently verifies 64 epoch files and recorded selection, and
independently recomputes selected-state objectives and complete-gallery
retention. It does not replay every epoch objective or the initial
gradient calibration. The development audit independently reconstructs
the aggregate statistics and continuation decisions. V10 effects come
from \path|results/practical_v10/official_development/joint_{encoder}_seed17/result.json|.

A workspace reset required restoring sources and issuing a replacement
LABCLIP-style contract. The restored fitter, runner and tests match their
original hashes. The replacement transparently records a new audit binding
and timestamp; the earlier contract hash remains historical. This recovery
does not add a recipe or outcome-dependent selection. The later evaluator
requires numerical replay of each selected LABCLIP-style objective and
identity/selected movement witnesses, plus verification of all 32 checkpoint
hashes and recorded minima; it does not claim numerical replay of all
32 audit objectives.

\section{Additional completed development intervals}
The source-pair intervals omitted from Table~\ref{tab:ci} are:
V8 ViT-B/32 $[-.930,.260]$, V8 RN50 $[-.698,.176]$,
V9 ViT-B/32 $[-1.457,1.506]$, V9 RN50 $[-.957,.860]$,
V10 ViT-B/32 $[-.713,.746]$, and V10 RN50 $[-.939,.891]$ percentage points.
All cross zero. The V9 and V10 original-joint composition calculations contain
15,630 declared triples and their source-pair calculations 5,340, each within
100 image clusters. These triple counts do not replace the image-cluster
sample size in uncertainty estimation.

\begin{thebibliography}{99}
\bibitem{clip} A. Radford et al.
Learning Transferable Visual Models From Natural Language Supervision.
ICML, 2021. \url{https://proceedings.mlr.press/v139/radford21a.html}.
\bibitem{pct} S. Yan, Y. Xiong, K. Kundu, S. Yang, S. Deng, M. Wang,
W. Xia, and S. Soatto.
Positive-Congruent Training: Towards Regression-Free Model Updates.
CVPR, 2021, pp. 14299--14308.
\url{https://openaccess.thecvf.com/content/CVPR2021/html/Yan_Positive-Congruent_Training_Towards_Regression-Free_Model_Updates_CVPR_2021_paper.html}.
\bibitem{bct} Y. Shen, Y. Xiong, W. Xia, and S. Soatto.
Towards Backward-Compatible Representation Learning. CVPR, 2020.
\url{https://arxiv.org/abs/2003.11942}.
\bibitem{rcat} D. Angioni, L. Demetrio, M. Pintor, L. Oneto, D. Anguita,
B. Biggio, and F. Roli.
Robustness-Congruent Adversarial Training for Secure Machine Learning Model
Updates. arXiv:2402.17390v2, 2025.
\url{https://arxiv.org/html/2402.17390v2}.
\bibitem{labclip} D. Koishigarina, A. Uselis, and S. J. Oh.
CLIP Behaves like a Bag-of-Words Model Cross-modally but not Uni-modally.
arXiv:2502.03566v3, 2026. \url{https://arxiv.org/html/2502.03566v3}.
\bibitem{clic} A. Peleg, N. D. Singh, and M. Hein.
Advancing Compositional Awareness in CLIP with Efficient Fine-Tuning.
NeurIPS, 2025. \url{https://arxiv.org/html/2505.24424v2}.
\bibitem{fsc} Y. Oh, J. W. Cho, D.-J. Kim, I. S. Kweon, and J. Kim.
Preserving Multi-Modal Capabilities of Pre-trained VLMs for Improving
Vision-Linguistic Compositionality. EMNLP, 2024.
\url{https://aclanthology.org/2024.emnlp-main.1062/}.
\bibitem{biclip} P. Mantini and S. K. Shah.
BiCLIP: Domain Canonicalization via Structured Geometric Transformation.
arXiv:2603.08942v2, 2026. \url{https://arxiv.org/abs/2603.08942v2}.
\bibitem{rankclip} Y. Zhang, Z. Zhao, Z. Chen, Z. Feng, Z. Ding, and Y. Sun.
RankCLIP: Ranking-Consistent Language-Image Pretraining.
arXiv:2404.09387v3, 2025. \url{https://arxiv.org/abs/2404.09387v3}.
\bibitem{inference} I. Miranda, A. Salaberria, E. Agirre, and G. Azkune.
Revisiting Compositionality in Dual-Encoder Vision-Language Models:
The Role of Inference. arXiv:2604.11496v2, 2026.
\url{https://arxiv.org/abs/2604.11496v2}.
\bibitem{evil} M. Kayser, O.-M. Camburu, L. Salewski, C. Emde, V. Do,
Z. Akata, and T. Lukasiewicz.
e-ViL: A Dataset and Benchmark for Natural Language Explanations in
Vision-Language Tasks. ICCV, 2021. \url{https://arxiv.org/abs/2105.03761v2}.
\bibitem{sugarcrepepp} S. H. Dumpala, A. Jaiswal, C. Sastry, E. Milios,
S. Oore, and H. Sajjad.
SugarCrepe++ Dataset: Vision-Language Model Sensitivity to Semantic and
Lexical Alterations. NeurIPS Datasets and Benchmarks, 2024.
\url{https://arxiv.org/abs/2406.11171}.
\bibitem{cose} D. Han, S. Moon, and S. Lee.
You CAN Teach an Old Model New Tricks: Domain Adaptation via
Complementary Subspace Expansion. Author project page, accessed
5 October 2026. \url{https://dhk1349.github.io/cose/}.
\end{thebibliography}
\end{document}
